\chapter{Zadatak 3a --- formulacija i klasa}
\label{ch:z3-formulacija}

\section{Što se ovdje traži}

Prevesti inženjerski opis --- ,,trošak je proporcionalan količini zemlje`` --- u
optimizacijski problem, i reći kojoj klasi taj problem pripada.

\begin{zadatak}{Zadatak 3, prvi podzadatak}
Formulirajte problem kao optimizacijski problem. Kojoj klasi optimizacijskih
problema pripada ovaj problem?
\end{zadatak}

\begin{predznanje}
\begin{itemize}[nosep,leftmargin=*]
  \item Što je konveksna funkcija --- \S\ref{sec:z1a-konveksan-skup}, Poglavlje~\ref{ch:z1-konveksnost}.
  \item Norma $\norm{\cdot}_1$ --- \S\ref{sec:vektori-matrice}.
  \item Značenje derivacije i druge derivacije --- \S\ref{sec:derivacije-oznake}.
\end{itemize}
\end{predznanje}

\section{Od zemljanih radova do integrala}
\label{sec:z3-model}

\begin{korak}[Izgradnja funkcije cilja, korak po korak]
\begin{enumerate}[leftmargin=*]
  \item \textbf{Što se plaća.} Trošak je proporcionalan količini zemlje koju
        treba otkopati \emph{ili} nasuti. Na mjestu $x$ razlika između nivelete i
        terena je $y(x) - h(x)$: pozitivna znači nasip, negativna iskop.
  \item \textbf{Zašto apsolutna vrijednost.} Oboje se plaća, pa u trošak ulazi
        $\abs{y(x) - h(x)}$. Bez apsolutne vrijednosti iskop bi ,,poništavao``
        nasip, pa bi optimalno rješenje bilo bilo koja krivulja koja u prosjeku
        prati teren --- što je fizikalno besmisleno.
  \item \textbf{Od visine do volumena.} Površina presjeka između dviju krivulja
        je $\int_0^L \abs{y - h}\dd x$. Uz jediničnu širinu trase, $w = 1$ m,
        volumen je brojčano jednak toj površini, izraženoj u $\mathrm{m}^2$.
  \item \textbf{Ograničenja.} Tri, izravno iz teksta: nagib
        $\abs{y'(x)} \leq b_1$, brzina promjene nagiba $\abs{y''(x)} \leq b_2$, i
        rubni uvjeti $y(0) = a$, $y(L) = b$.
\end{enumerate}
\end{korak}

\begin{formula}{Kontinuirana formulacija}
\begin{equation}
\min_{y(\cdot)} \ J[y] = \int_0^L \bigl| y(x) - h(x) \bigr| \dd x
\label{eq:z3-kont}
\end{equation}
uz ograničenja
\begin{equation}
\abs{y'(x)} \leq b_1, \qquad \abs{y''(x)} \leq b_2, \qquad y(0) = a, \quad y(L) = b .
\label{eq:z3-kont-ogr}
\end{equation}
\end{formula}

Zadani podaci:
\begin{equation}
L = 10\ \mathrm{m}, \qquad
h(x) = 0{,}0053x^4 - 0{,}095x^3 + 0{,}48x^2 - 0{,}3x + 1,
\end{equation}
\begin{equation}
a = h(0) = 1\ \mathrm{m}, \qquad b = h(L) = 4\ \mathrm{m},
\end{equation}
\begin{equation}
\text{slučaj A:}\ b_1 = 0{,}5,\ b_2 = 0{,}5\ \mathrm{m}^{-1};
\qquad
\text{slučaj B:}\ b_1 = 0{,}5,\ b_2 = 0{,}2\ \mathrm{m}^{-1}.
\end{equation}

\begin{korak}[Provjera rubnih uvjeta]
Zadatak tvrdi $a = h(0) = 1$ i $b = h(L) = 4$. Provjerimo:
\begin{align}
h(0) &= 1, \\
h(10) &= 0{,}0053 \cdot 10^4 - 0{,}095 \cdot 10^3 + 0{,}48 \cdot 10^2 - 0{,}3 \cdot 10 + 1 \\
      &= 53 - 95 + 48 - 3 + 1 = 4 . \quad\checkmark
\end{align}
Rubni uvjeti, dakle, traže da niveleta \emph{počinje i završava na terenu} ---
ondje nema ni iskopa ni nasipa. To je važno za razumijevanje slike: odstupanje
mora biti nula na oba kraja.
\end{korak}

\section{Što $b_1$ i $b_2$ fizikalno znače}
\label{sec:z3-fizikalno}

\begin{pojam}{Nagib i brzina promjene nagiba}
\begin{itemize}[leftmargin=*]
  \item $b_1$ ograničava \textbf{nagib} nivelete, $y'(x)$. Bezdimenzijski je
        (m/m); $b_1 = 0{,}5$ znači najviše $50\%$ uspona, odnosno oko $26{,}6^\circ$.
        To je uvjet izvedivosti --- vozilo mora moći savladati uspon.
  \item $b_2$ ograničava \textbf{brzinu promjene nagiba}, $y''(x)$. Jedinica je
        $\mathrm{m}^{-1}$. Pri malim nagibima ovo je aproksimacija vertikalne
        \emph{zakrivljenosti} trase, jer je
        $\kappa = y''/(1 + y'^2)^{3/2} \approx y''$ za mali $y'$. To je uvjet
        udobnosti --- zakrivljenost izravno određuje vertikalno ubrzanje koje
        putnik osjeća pri danoj brzini vožnje.
\end{itemize}
\end{pojam}

\begin{zamka}[Zamka: $y''$ nije zakrivljenost, nego njezina aproksimacija]
Prava zakrivljenost krivulje je $\kappa = y''/(1+y'^2)^{3/2}$. Uz $y' = 0{,}5$
nazivnik iznosi $(1{,}25)^{3/2} \approx 1{,}40$, pa $y''$ precjenjuje
zakrivljenost za oko $40\%$ na najstrmijim dionicama. Zadatak izričito traži
$y''$, pa se tako i radi --- ali ako te se pita ,,je li to zakrivljenost``,
točan odgovor je: aproksimacija koja vrijedi za male nagibe, a ovdje nagibi nisu
zanemarivi.
\end{zamka}

\section{Kojoj klasi problem pripada}
\label{sec:z3-klasa}

Odgovor je \textbf{linearni program}, ali tek nakon dva koraka: diskretizacije i
uklanjanja apsolutne vrijednosti. Prije toga problem nije ni konačnodimenzionalan
(nepoznanica je \emph{funkcija} $y(\cdot)$), a funkcija cilja nije glatka.

\begin{izvor}{Predavanje 5 --- LP i QP, slajd 4/34 (PDF str.~4)}
\slajd[0.86]{p05-lp}

\textbf{Transkripcija.} \emph{LP problem.} Za zadane matrice $G \in \R^{p\times n}$,
$A \in \R^{q\times n}$ i vektore $c \in \R^n$, $h \in \R^p$, $b \in \R^q$,
problem linearnog programiranja je optimizacijski problem oblika
\[
\begin{aligned}
\text{minimiziraj} \quad & c\trans x \\
\text{uz ograničenja} \quad & Gx \leq h \\
& Ax = b
\end{aligned}
\]
gdje je $x \in \R^n$ vektor projektnih varijabli.
\begin{itemize}[nosep,leftmargin=*]
  \item Funkcija cilja je linearna u projektnim varijablama.
  \item Ograničenja su linearna (afina).
  \item LP problemi i sa preko $10\,000$ varijabli se mogu učinkovito riješiti.
  \item Široko dostupni računalni programi za rješavanje LP problema.
  \item \emph{,,Rješavanje LP problema je tehnologija!``}
\end{itemize}

\textbf{Zašto to citiramo:} zadnja rečenica objašnjava \emph{zašto} se isplati
mukotrpno svoditi problem na LP. Naš problem s $N = 100$ ima $202$ varijable i
$600$ nejednakosti --- za LP to je sitno.
\end{izvor}

\begin{izvor}{Vježbe 4 --- slajd 4/40 (PDF str.~4)}
\slajd[0.86]{v4-aproksimacija}

\textbf{Transkripcija.} \emph{Problem aproksimacije.} Za zadane
$A \in \R^{m\times n}$, gdje je $m > n$, i $b \in \R^m$:
\[
\min_x \ \norm{Ax - b}
\]
gdje je $\norm{\cdot}$ zadana norma. Interpretacije:
\begin{itemize}[nosep,leftmargin=*]
  \item \emph{aproksimacija}: vektor $A\xopt$ je najbolja aproksimacija vektora
        $b$ dobivena linearnom kombinacijom stupaca od $A$;
  \item \emph{geometrijska}: $A\xopt$ je točka u slici od $A$ koja je najbliža
        $b$-u (u promatranoj normi);
  \item \emph{estimacija}: za linearan model $y = Ax + v$ mjerimo $y$ a želimo
        odrediti $x$, dok je $v$ šum;
  \item \emph{optimalan dizajn}: $x$ je vektor ulaznih varijabli, $Ax$ rezultat;
        za željeni rezultat $b$, $\xopt$ je najbolji izbor.
\end{itemize}

\textbf{Naš zadatak je četvrta interpretacija:} $y$ su projektne varijable
(visine nivelete), $h$ je ,,željeni rezultat`` (teren), a tražimo najbolju
aproksimaciju terena u normi $\ell_1$ --- ali s dodatnim ograničenjima na
glatkoću, kakvih u čistom problemu aproksimacije nema.
\end{izvor}

\begin{formula}{Klasa problema}
Nakon diskretizacije i uklanjanja apsolutne vrijednosti funkcija cilja i sva
ograničenja su linearna u optimizacijskim varijablama, pa problem pripada klasi
\textbf{linearnog programiranja (LP)}.

Preciznije, riječ je o \textbf{$\ell_1$-aproksimaciji s dodatnim linearnim
ograničenjima}: minimizira se $\norm{y - h}_1$ uz linearna ograničenja na $D_1 y$
i $D_2 y$ i rubne uvjete.
\end{formula}

\begin{izgovor}
,,Trošak je proporcionalan količini premještene zemlje, dakle površini između
nivelete i terena. Kako se i iskop i nasip plaćaju, u kriterij ulazi apsolutna
vrijednost razlike, pa je kontinuirani problem: minimiziraj integral apsolutne
vrijednosti ipsilon minus ha, po iks od nule do el, uz ograničenja na prvu i
drugu derivaciju i uz zadane rubne uvjete.

Klasa: nakon diskretizacije i uklanjanja apsolutne vrijednosti pomoćnim
varijablama, sve postaje linearno --- dakle linearni program. Preciznije, to je
$\ell_1$-aproksimacija s linearnim ograničenjima, u obliku koji smo radili na
Vježbama 4.

Naglasio bih da prije ta dva koraka problem nije LP: nepoznanica je funkcija, ne
vektor, a apsolutna vrijednost nije glatka. Konveksan jest od samog početka, jer
je apsolutna vrijednost konveksna, a ograničenja afina.``
\end{izgovor}
